\documentclass[11pt]{article}
\usepackage[margin=1in]{geometry}
\usepackage{amsmath,amssymb,bm,braket}
\usepackage{hyperref}
\begin{document}
\section*{Problem setup}
A magic trapping wavelength gives equal AC Stark shifts for two states of
an imaging transition. Near-magic denotes a reported small residual differential
shift. Consider imaging on the $^1S_0$--$^3P_1$ line of $^{171}$Yb and $^{174}$Yb.
The trapping wavelength and the imaging-light wavelength are distinct.
\section*{Main problem}
Within 400--600 nm, identify the following reported or explicitly inferred
imaging conditions, treating their experimental configurations separately:
\begin{enumerate}
\item The $^{174}$Yb $m_J=0\to m'_J=0$ near-magic condition for linearly polarized
tweezers with no applied magnetic field, quantized along the trapping
polarization, in Saskin et al., \url{https://arxiv.org/abs/1810.10517}.
\item The $^{171}$Yb $F=1/2\to F'=3/2$, $|m'_F|=1/2$ condition in the
configuration of Ma et al., \url{https://arxiv.org/abs/2112.06799}.
\item The $^{171}$Yb stretched $|m_F|=1/2\to |m'_F|=3/2$ imaging condition
with trap polarization perpendicular to the magnetic field in Norcia et al.,
\url{https://arxiv.org/abs/2305.19119}.
\item The corresponding $^{174}$Yb $m_J=0\to m'_J=\pm1$ estimate inferred
from the preceding stretched-state electronic polarizability. For this
inference, neglect isotope/hyperfine changes to the electronic polarizability
and retain the same electronic-state and trap-polarization geometry.
\end{enumerate}
For each case give the approximate trapping wavelength, the addressed states,
and the imaging transition type ($\pi$ for $\Delta m=0$, $\sigma$ for
$\Delta m=\pm1$). Distinguish measured magic, measured near-magic and inferred
values, and do not imply a common geometry or more precision than the sources
support. This is a literature-based identification of these specified cases,
not a first-principles polarizability calculation or an exhaustive search for
all crossings in this interval.
\end{document}
